% ==========================================================================
% interface.tex -- working notes: SBM for jump-coefficient (interface)
% problems WITHOUT a mixed formulation. Derivation of a consistent shifted
% interface coupling, diagnosis of why the naive form fails (already in 1D),
% and numerical validation. Reproduce all numbers: python3 interface_1d.py
% % ==========================================================================
% interface.tex -- working notes: SBM for jump-coefficient (interface)
% problems WITHOUT a mixed formulation. Derivation of a consistent shifted
% interface coupling, diagnosis of why the naive form fails (already in 1D),
% and numerical validation. Reproduce all numbers: python3 interface_1d.py
% % ==========================================================================
% interface.tex -- working notes: SBM for jump-coefficient (interface)
% problems WITHOUT a mixed formulation. Derivation of a consistent shifted
% interface coupling, diagnosis of why the naive form fails (already in 1D),
% and numerical validation. Reproduce all numbers: python3 interface_1d.py
% % ==========================================================================
% interface.tex -- working notes: SBM for jump-coefficient (interface)
% problems WITHOUT a mixed formulation. Derivation of a consistent shifted
% interface coupling, diagnosis of why the naive form fails (already in 1D),
% and numerical validation. Reproduce all numbers: python3 interface_1d.py
% \input{interface.tex} at the end of the paper to compile.
% ==========================================================================

\section*{Appendix (working notes): a consistent shifted formulation for interface problems}

\subsection*{I.1 Problem and setup}

Model problem on $(0,1)$ (all statements extend verbatim to $d>1$):
\[
  -(\kappa u')' = f, \qquad
  \kappa = \begin{cases} \kappa^- & x < x_\Gamma \\ \kappa^+ & x > x_\Gamma \end{cases},
  \qquad
  [u]_{x_\Gamma} = 0, \quad [\kappa u']_{x_\Gamma} = 0,
\]
with the interface $x_\Gamma$ \emph{not} aligned with the mesh. Shifted
treatment: the surrogate interface is the nearest mesh vertex $\tilde x$,
shift $d = x_\Gamma - \tilde x$ (so $|d| \le h/2$ by construction --- note
this is automatically milder than the Dirichlet SBM staircase, where $s$
reaches $\sqrt 2$). The degree of freedom at $\tilde x$ is duplicated: two
one-sided fields $u_h^\pm$, each assembled with its \emph{own} constant
$\kappa^\pm$ on its own surrogate subdomain.

Interpretation that makes everything work: $u_h^\pm$ do \emph{not}
approximate the physical solution near the interface; they approximate the
smooth \emph{extensions} $\tilde u^\pm$ of the two physical pieces, each
solving its own $\kappa^\pm$-equation past $x_\Gamma$. In the strip between
$\tilde x$ and $x_\Gamma$, the physical solution is recovered in
post-processing from the field of the side that physically owns the strip
(evaluated by extension). The interface conditions are imposed at the
\emph{true} interface via Taylor/polynomial extension
$\mathcal E^\pm w (x_\Gamma) = w(\text{eval.\ of the adjacent element's
polynomial at } x_\Gamma)$, exactly as in the Dirichlet SBM. One side
evaluates by extrapolation (distance $|d|$), the other by interior
evaluation. No mixed formulation, no extra flux variable.

\subsection*{I.2 Why the natural symmetric form fails (diagnosis of the 1D failure)}

Notation: plain jump at the surrogate vertex $[v]_{\tilde x} = v^-(\tilde
x) - v^+(\tilde x)$; extended jump at the true interface $[\mathcal E
v]_\Gamma = \mathcal E^- v^- (x_\Gamma) - \mathcal E^+ v^+(x_\Gamma)$;
weighted average of extended fluxes $\{\kappa \mathcal E u'\}_\Gamma =
\beta^- \kappa^- \mathcal E^- u'\,(x_\Gamma) + \beta^+ \kappa^+ \mathcal
E^+ u'(x_\Gamma)$ with the harmonic weights $\beta^\mp =
\kappa^\pm/(\kappa^-{+}\kappa^+)$.

The trap is a mismatch of \emph{locations}: integration by parts of the
exact solution over the two surrogate subdomains produces flux terms
\emph{at $\tilde x$} acting on the \emph{plain} jump $[v]_{\tilde x}$,
whereas the physical conditions hold \emph{at $x_\Gamma$} in the
\emph{extended} variables. The symmetric-looking Nitsche form that pairs
the flux with $[\mathcal E v]_\Gamma$,
\[
  \cdots - \{\kappa \mathcal E u'\}_\Gamma\, [\mathcal E v]_\Gamma - \cdots
  \qquad \text{(``naive'')},
\]
leaves, for the exact broken line with flux $F$ (take $f=0$, $p=1$), the
residual
\[
  F\big( [v]_{\tilde x} - [\mathcal E v]_\Gamma \big)
  = -F\, d\, [v']_{\tilde x} \;\neq\; 0 ,
\]
an $O(d) = O(h)$ \emph{consistency} error concentrated at the interface:
the method converges to the wrong solution, and already the piecewise
linear exact solution is not reproduced. This is precisely the observed 1D
failure. (Verified numerically: the naive form gives errors
$3\cdot10^{-2}$--$2\cdot10^{-1}$ on the broken line, non-decreasing in
$p$.)

\subsection*{I.3 The consistent form}

The two interface conditions play different roles and must live at
different locations:

\begin{equation}
  a(u,v) \;=\; \sum_\pm \int_{\tilde\Omega^\pm} \kappa\, u' v'
  \;-\; \{\kappa \mathcal E u'\}_\Gamma\, \underbrace{[v]_{\tilde x}}_{\text{plain, at }\tilde x}
  \;-\; \{\kappa \mathcal E v'\}_\Gamma\, \underbrace{[\mathcal E u]_\Gamma}_{\text{extended, at }\Gamma}
  \;+\; \sigma\, [\mathcal E u]_\Gamma [\mathcal E v]_\Gamma ,
  \label{eq:sim_form}
\end{equation}
with $\sigma = \gamma\, \kappa_{\mathrm{hm}}\, p^2/h$,
$\kappa_{\mathrm{hm}} = 2\kappa^-\kappa^+/(\kappa^-{+}\kappa^+)$, and the
right-hand side carrying a \emph{shifted-source correction}:
\begin{equation}
  L(v) \;=\; \sum_\pm \int_{\tilde\Omega^\pm} f v
  \;+\; \Big( \int_{\tilde x}^{x_\Gamma} f \, dx \Big)\, [v]_{\tilde x} .
  \label{eq:sim_rhs}
\end{equation}

Rationale, term by term:
\begin{itemize}
  \item \textbf{First flux term at $\tilde x$, extended flux from
        $\Gamma$:} integration by parts puts the flux terms at $\tilde x$;
        the physically known quantity (continuous flux) lives at
        $x_\Gamma$. For $f = 0$ the two coincide on each side (the flux of
        the extension is constant in 1D); the extended flux average is
        exactly what cancels the IBP terms. This mirrors the structure of
        the Dirichlet SBM form in the paper, where the first flux term
        also carries the \emph{plain} test trace and only the adjoint
        term carries $\mathcal E$.
  \item \textbf{Adjoint term and penalty on $[\mathcal E u]_\Gamma$:} the
        value-continuity condition holds at $x_\Gamma$; both terms vanish
        on the exact solution, so any consistent-looking location would
        do --- placing them at $\Gamma$ in extended variables enforces
        the correct condition to full order.
  \item \textbf{Shifted source \eqref{eq:sim_rhs}:} with $f \ne 0$ the
        one-sided fluxes at $\tilde x$ and at $x_\Gamma$ differ by the
        strip integral, $\kappa u'(\tilde x) = \kappa u'(x_\Gamma) +
        \int_{\tilde x}^{x_\Gamma} f$; for $f$ continuous this shift is
        the \emph{same on both sides}, hence it multiplies
        $[v]_{\tilde x}$ and moves to the right-hand side. Omitting it
        costs $O(h)$ consistency (observed: convergence stalls). For $f$
        discontinuous across $\Gamma$ the two shifts differ and an
        additional $\{v\}$-weighted data term appears --- same derivation.
\end{itemize}

\subsection*{I.4 Numerical validation (see \texttt{interface\_1d.py})}

$\kappa^- = 1$, $\kappa^+ = 10$, $x_\Gamma = 0.523$ (generic, off-node),
$\mathbb P_p$ elements, $n = 8$ elements unless stated:

\begin{center}
\begin{tabular}{lccc}
\hline
 & $p=1$ & $p=2$ & $p=3$ \\
\hline
broken line ($f=0$), consistent form & $8\cdot10^{-15}$ & $5\cdot10^{-14}$ & $8\cdot10^{-14}$ \\
broken line, naive form              & $3\cdot10^{-2}$ & $8\cdot10^{-2}$ & $2\cdot10^{-1}$ \\
$f=1$ (piecewise quadratic), $p\ge2$ &  --- & $7\cdot10^{-14}$ & $6\cdot10^{-14}$ \\
\hline
\end{tabular}
\end{center}

The consistent form reproduces every exact solution lying in the discrete
space to machine precision (broken line for all $p$; the piecewise
quadratic for $p \ge 2$, which certifies the shifted-source term), while
$p=1$ converges at the expected rate. Robustness in the jump: with
$\kappa^+/\kappa^- = 10^6$ the broken line is reproduced to $10^{-8}$
(conditioning-limited), thanks to the harmonic weights
$\beta^\mp \propto \kappa^\pm$ and the harmonic-mean penalty --- the
standard cure from unfitted Nitsche interface methods, which carries over
unchanged.

\subsection*{I.5 Remarks and outlook to 2D/3D}

\begin{itemize}
  \item \textbf{Relation to the literature.} A shifted interface method
        of this flavor exists (Li \& Scovazzi and follow-ups); these
        notes' added value is the explicit identification of (a) the
        location mismatch $[v]_{\tilde x}$ vs.\ $[\mathcal E v]_\Gamma$ as
        the reason naive derivations fail already in 1D, and (b) the
        shifted-source correction \eqref{eq:sim_rhs}, which is easy to
        miss (it vanishes in every $f=0$ sanity test) yet caps accuracy
        at first order when omitted.
  \item \textbf{No mixed formulation is needed.} Flux continuity is
        enforced through the extended-flux average exactly as value
        continuity is enforced through the extended jump; both are
        Nitsche-type terms on the primal variable.
  \item \textbf{In 2D} the surrogate interface is the staircase of faces
        nearest to $\Gamma$, the strip cells are assigned to one side
        with that side's $\kappa$, and all findings of the definiteness
        notes (theory.tex) apply with two changes in our favor: the
        shift satisfies $|d| \le h/2$ per face by nearest-face selection
        (vs.\ $\sqrt2 h$ for the Dirichlet staircase pockets), and both
        traces are one-sided extrapolations over that reduced distance.
        The same Chebyshev certificates should be run for the coupled
        two-cell block (two fields, duplicated interface DoFs); the
        chamfer and the PDE-substituted extension
        ($\partial_n^2 u = -(f + \nabla\kappa\cdot\nabla u)/\kappa
        - \Delta_{\mathrm{tan}} u$ per side, $\kappa$ constant per side)
        transfer verbatim.
  \item \textbf{Smoother.} The duplicated interface DoFs must belong to a
        common patch (the vertex patch at $\tilde x$ contains both
        one-sided fields by construction), so the full-residual shy patch
        machinery extends without structural change.
  \item \textbf{Open point.} The first flux term in \eqref{eq:sim_form}
        is non-symmetric by construction (plain vs.\ extended test
        traces); the effect on the definiteness certificates at
        $p \ge 3$, and whether the interface analogue prefers the
        PDE-substituted extension from the start, should be checked with
        the same single-block methodology before a 2D implementation.
\end{itemize}
 at the end of the paper to compile.
% ==========================================================================

\section*{Appendix (working notes): a consistent shifted formulation for interface problems}

\subsection*{I.1 Problem and setup}

Model problem on $(0,1)$ (all statements extend verbatim to $d>1$):
\[
  -(\kappa u')' = f, \qquad
  \kappa = \begin{cases} \kappa^- & x < x_\Gamma \\ \kappa^+ & x > x_\Gamma \end{cases},
  \qquad
  [u]_{x_\Gamma} = 0, \quad [\kappa u']_{x_\Gamma} = 0,
\]
with the interface $x_\Gamma$ \emph{not} aligned with the mesh. Shifted
treatment: the surrogate interface is the nearest mesh vertex $\tilde x$,
shift $d = x_\Gamma - \tilde x$ (so $|d| \le h/2$ by construction --- note
this is automatically milder than the Dirichlet SBM staircase, where $s$
reaches $\sqrt 2$). The degree of freedom at $\tilde x$ is duplicated: two
one-sided fields $u_h^\pm$, each assembled with its \emph{own} constant
$\kappa^\pm$ on its own surrogate subdomain.

Interpretation that makes everything work: $u_h^\pm$ do \emph{not}
approximate the physical solution near the interface; they approximate the
smooth \emph{extensions} $\tilde u^\pm$ of the two physical pieces, each
solving its own $\kappa^\pm$-equation past $x_\Gamma$. In the strip between
$\tilde x$ and $x_\Gamma$, the physical solution is recovered in
post-processing from the field of the side that physically owns the strip
(evaluated by extension). The interface conditions are imposed at the
\emph{true} interface via Taylor/polynomial extension
$\mathcal E^\pm w (x_\Gamma) = w(\text{eval.\ of the adjacent element's
polynomial at } x_\Gamma)$, exactly as in the Dirichlet SBM. One side
evaluates by extrapolation (distance $|d|$), the other by interior
evaluation. No mixed formulation, no extra flux variable.

\subsection*{I.2 Why the natural symmetric form fails (diagnosis of the 1D failure)}

Notation: plain jump at the surrogate vertex $[v]_{\tilde x} = v^-(\tilde
x) - v^+(\tilde x)$; extended jump at the true interface $[\mathcal E
v]_\Gamma = \mathcal E^- v^- (x_\Gamma) - \mathcal E^+ v^+(x_\Gamma)$;
weighted average of extended fluxes $\{\kappa \mathcal E u'\}_\Gamma =
\beta^- \kappa^- \mathcal E^- u'\,(x_\Gamma) + \beta^+ \kappa^+ \mathcal
E^+ u'(x_\Gamma)$ with the harmonic weights $\beta^\mp =
\kappa^\pm/(\kappa^-{+}\kappa^+)$.

The trap is a mismatch of \emph{locations}: integration by parts of the
exact solution over the two surrogate subdomains produces flux terms
\emph{at $\tilde x$} acting on the \emph{plain} jump $[v]_{\tilde x}$,
whereas the physical conditions hold \emph{at $x_\Gamma$} in the
\emph{extended} variables. The symmetric-looking Nitsche form that pairs
the flux with $[\mathcal E v]_\Gamma$,
\[
  \cdots - \{\kappa \mathcal E u'\}_\Gamma\, [\mathcal E v]_\Gamma - \cdots
  \qquad \text{(``naive'')},
\]
leaves, for the exact broken line with flux $F$ (take $f=0$, $p=1$), the
residual
\[
  F\big( [v]_{\tilde x} - [\mathcal E v]_\Gamma \big)
  = -F\, d\, [v']_{\tilde x} \;\neq\; 0 ,
\]
an $O(d) = O(h)$ \emph{consistency} error concentrated at the interface:
the method converges to the wrong solution, and already the piecewise
linear exact solution is not reproduced. This is precisely the observed 1D
failure. (Verified numerically: the naive form gives errors
$3\cdot10^{-2}$--$2\cdot10^{-1}$ on the broken line, non-decreasing in
$p$.)

\subsection*{I.3 The consistent form}

The two interface conditions play different roles and must live at
different locations:

\begin{equation}
  a(u,v) \;=\; \sum_\pm \int_{\tilde\Omega^\pm} \kappa\, u' v'
  \;-\; \{\kappa \mathcal E u'\}_\Gamma\, \underbrace{[v]_{\tilde x}}_{\text{plain, at }\tilde x}
  \;-\; \{\kappa \mathcal E v'\}_\Gamma\, \underbrace{[\mathcal E u]_\Gamma}_{\text{extended, at }\Gamma}
  \;+\; \sigma\, [\mathcal E u]_\Gamma [\mathcal E v]_\Gamma ,
  \label{eq:sim_form}
\end{equation}
with $\sigma = \gamma\, \kappa_{\mathrm{hm}}\, p^2/h$,
$\kappa_{\mathrm{hm}} = 2\kappa^-\kappa^+/(\kappa^-{+}\kappa^+)$, and the
right-hand side carrying a \emph{shifted-source correction}:
\begin{equation}
  L(v) \;=\; \sum_\pm \int_{\tilde\Omega^\pm} f v
  \;+\; \Big( \int_{\tilde x}^{x_\Gamma} f \, dx \Big)\, [v]_{\tilde x} .
  \label{eq:sim_rhs}
\end{equation}

Rationale, term by term:
\begin{itemize}
  \item \textbf{First flux term at $\tilde x$, extended flux from
        $\Gamma$:} integration by parts puts the flux terms at $\tilde x$;
        the physically known quantity (continuous flux) lives at
        $x_\Gamma$. For $f = 0$ the two coincide on each side (the flux of
        the extension is constant in 1D); the extended flux average is
        exactly what cancels the IBP terms. This mirrors the structure of
        the Dirichlet SBM form in the paper, where the first flux term
        also carries the \emph{plain} test trace and only the adjoint
        term carries $\mathcal E$.
  \item \textbf{Adjoint term and penalty on $[\mathcal E u]_\Gamma$:} the
        value-continuity condition holds at $x_\Gamma$; both terms vanish
        on the exact solution, so any consistent-looking location would
        do --- placing them at $\Gamma$ in extended variables enforces
        the correct condition to full order.
  \item \textbf{Shifted source \eqref{eq:sim_rhs}:} with $f \ne 0$ the
        one-sided fluxes at $\tilde x$ and at $x_\Gamma$ differ by the
        strip integral, $\kappa u'(\tilde x) = \kappa u'(x_\Gamma) +
        \int_{\tilde x}^{x_\Gamma} f$; for $f$ continuous this shift is
        the \emph{same on both sides}, hence it multiplies
        $[v]_{\tilde x}$ and moves to the right-hand side. Omitting it
        costs $O(h)$ consistency (observed: convergence stalls). For $f$
        discontinuous across $\Gamma$ the two shifts differ and an
        additional $\{v\}$-weighted data term appears --- same derivation.
\end{itemize}

\subsection*{I.4 Numerical validation (see \texttt{interface\_1d.py})}

$\kappa^- = 1$, $\kappa^+ = 10$, $x_\Gamma = 0.523$ (generic, off-node),
$\mathbb P_p$ elements, $n = 8$ elements unless stated:

\begin{center}
\begin{tabular}{lccc}
\hline
 & $p=1$ & $p=2$ & $p=3$ \\
\hline
broken line ($f=0$), consistent form & $8\cdot10^{-15}$ & $5\cdot10^{-14}$ & $8\cdot10^{-14}$ \\
broken line, naive form              & $3\cdot10^{-2}$ & $8\cdot10^{-2}$ & $2\cdot10^{-1}$ \\
$f=1$ (piecewise quadratic), $p\ge2$ &  --- & $7\cdot10^{-14}$ & $6\cdot10^{-14}$ \\
\hline
\end{tabular}
\end{center}

The consistent form reproduces every exact solution lying in the discrete
space to machine precision (broken line for all $p$; the piecewise
quadratic for $p \ge 2$, which certifies the shifted-source term), while
$p=1$ converges at the expected rate. Robustness in the jump: with
$\kappa^+/\kappa^- = 10^6$ the broken line is reproduced to $10^{-8}$
(conditioning-limited), thanks to the harmonic weights
$\beta^\mp \propto \kappa^\pm$ and the harmonic-mean penalty --- the
standard cure from unfitted Nitsche interface methods, which carries over
unchanged.

\subsection*{I.5 Remarks and outlook to 2D/3D}

\begin{itemize}
  \item \textbf{Relation to the literature.} A shifted interface method
        of this flavor exists (Li \& Scovazzi and follow-ups); these
        notes' added value is the explicit identification of (a) the
        location mismatch $[v]_{\tilde x}$ vs.\ $[\mathcal E v]_\Gamma$ as
        the reason naive derivations fail already in 1D, and (b) the
        shifted-source correction \eqref{eq:sim_rhs}, which is easy to
        miss (it vanishes in every $f=0$ sanity test) yet caps accuracy
        at first order when omitted.
  \item \textbf{No mixed formulation is needed.} Flux continuity is
        enforced through the extended-flux average exactly as value
        continuity is enforced through the extended jump; both are
        Nitsche-type terms on the primal variable.
  \item \textbf{In 2D} the surrogate interface is the staircase of faces
        nearest to $\Gamma$, the strip cells are assigned to one side
        with that side's $\kappa$, and all findings of the definiteness
        notes (theory.tex) apply with two changes in our favor: the
        shift satisfies $|d| \le h/2$ per face by nearest-face selection
        (vs.\ $\sqrt2 h$ for the Dirichlet staircase pockets), and both
        traces are one-sided extrapolations over that reduced distance.
        The same Chebyshev certificates should be run for the coupled
        two-cell block (two fields, duplicated interface DoFs); the
        chamfer and the PDE-substituted extension
        ($\partial_n^2 u = -(f + \nabla\kappa\cdot\nabla u)/\kappa
        - \Delta_{\mathrm{tan}} u$ per side, $\kappa$ constant per side)
        transfer verbatim.
  \item \textbf{Smoother.} The duplicated interface DoFs must belong to a
        common patch (the vertex patch at $\tilde x$ contains both
        one-sided fields by construction), so the full-residual shy patch
        machinery extends without structural change.
  \item \textbf{Open point.} The first flux term in \eqref{eq:sim_form}
        is non-symmetric by construction (plain vs.\ extended test
        traces); the effect on the definiteness certificates at
        $p \ge 3$, and whether the interface analogue prefers the
        PDE-substituted extension from the start, should be checked with
        the same single-block methodology before a 2D implementation.
\end{itemize}
 at the end of the paper to compile.
% ==========================================================================

\section*{Appendix (working notes): a consistent shifted formulation for interface problems}

\subsection*{I.1 Problem and setup}

Model problem on $(0,1)$ (all statements extend verbatim to $d>1$):
\[
  -(\kappa u')' = f, \qquad
  \kappa = \begin{cases} \kappa^- & x < x_\Gamma \\ \kappa^+ & x > x_\Gamma \end{cases},
  \qquad
  [u]_{x_\Gamma} = 0, \quad [\kappa u']_{x_\Gamma} = 0,
\]
with the interface $x_\Gamma$ \emph{not} aligned with the mesh. Shifted
treatment: the surrogate interface is the nearest mesh vertex $\tilde x$,
shift $d = x_\Gamma - \tilde x$ (so $|d| \le h/2$ by construction --- note
this is automatically milder than the Dirichlet SBM staircase, where $s$
reaches $\sqrt 2$). The degree of freedom at $\tilde x$ is duplicated: two
one-sided fields $u_h^\pm$, each assembled with its \emph{own} constant
$\kappa^\pm$ on its own surrogate subdomain.

Interpretation that makes everything work: $u_h^\pm$ do \emph{not}
approximate the physical solution near the interface; they approximate the
smooth \emph{extensions} $\tilde u^\pm$ of the two physical pieces, each
solving its own $\kappa^\pm$-equation past $x_\Gamma$. In the strip between
$\tilde x$ and $x_\Gamma$, the physical solution is recovered in
post-processing from the field of the side that physically owns the strip
(evaluated by extension). The interface conditions are imposed at the
\emph{true} interface via Taylor/polynomial extension
$\mathcal E^\pm w (x_\Gamma) = w(\text{eval.\ of the adjacent element's
polynomial at } x_\Gamma)$, exactly as in the Dirichlet SBM. One side
evaluates by extrapolation (distance $|d|$), the other by interior
evaluation. No mixed formulation, no extra flux variable.

\subsection*{I.2 Why the natural symmetric form fails (diagnosis of the 1D failure)}

Notation: plain jump at the surrogate vertex $[v]_{\tilde x} = v^-(\tilde
x) - v^+(\tilde x)$; extended jump at the true interface $[\mathcal E
v]_\Gamma = \mathcal E^- v^- (x_\Gamma) - \mathcal E^+ v^+(x_\Gamma)$;
weighted average of extended fluxes $\{\kappa \mathcal E u'\}_\Gamma =
\beta^- \kappa^- \mathcal E^- u'\,(x_\Gamma) + \beta^+ \kappa^+ \mathcal
E^+ u'(x_\Gamma)$ with the harmonic weights $\beta^\mp =
\kappa^\pm/(\kappa^-{+}\kappa^+)$.

The trap is a mismatch of \emph{locations}: integration by parts of the
exact solution over the two surrogate subdomains produces flux terms
\emph{at $\tilde x$} acting on the \emph{plain} jump $[v]_{\tilde x}$,
whereas the physical conditions hold \emph{at $x_\Gamma$} in the
\emph{extended} variables. The symmetric-looking Nitsche form that pairs
the flux with $[\mathcal E v]_\Gamma$,
\[
  \cdots - \{\kappa \mathcal E u'\}_\Gamma\, [\mathcal E v]_\Gamma - \cdots
  \qquad \text{(``naive'')},
\]
leaves, for the exact broken line with flux $F$ (take $f=0$, $p=1$), the
residual
\[
  F\big( [v]_{\tilde x} - [\mathcal E v]_\Gamma \big)
  = -F\, d\, [v']_{\tilde x} \;\neq\; 0 ,
\]
an $O(d) = O(h)$ \emph{consistency} error concentrated at the interface:
the method converges to the wrong solution, and already the piecewise
linear exact solution is not reproduced. This is precisely the observed 1D
failure. (Verified numerically: the naive form gives errors
$3\cdot10^{-2}$--$2\cdot10^{-1}$ on the broken line, non-decreasing in
$p$.)

\subsection*{I.3 The consistent form}

The two interface conditions play different roles and must live at
different locations:

\begin{equation}
  a(u,v) \;=\; \sum_\pm \int_{\tilde\Omega^\pm} \kappa\, u' v'
  \;-\; \{\kappa \mathcal E u'\}_\Gamma\, \underbrace{[v]_{\tilde x}}_{\text{plain, at }\tilde x}
  \;-\; \{\kappa \mathcal E v'\}_\Gamma\, \underbrace{[\mathcal E u]_\Gamma}_{\text{extended, at }\Gamma}
  \;+\; \sigma\, [\mathcal E u]_\Gamma [\mathcal E v]_\Gamma ,
  \label{eq:sim_form}
\end{equation}
with $\sigma = \gamma\, \kappa_{\mathrm{hm}}\, p^2/h$,
$\kappa_{\mathrm{hm}} = 2\kappa^-\kappa^+/(\kappa^-{+}\kappa^+)$, and the
right-hand side carrying a \emph{shifted-source correction}:
\begin{equation}
  L(v) \;=\; \sum_\pm \int_{\tilde\Omega^\pm} f v
  \;+\; \Big( \int_{\tilde x}^{x_\Gamma} f \, dx \Big)\, [v]_{\tilde x} .
  \label{eq:sim_rhs}
\end{equation}

Rationale, term by term:
\begin{itemize}
  \item \textbf{First flux term at $\tilde x$, extended flux from
        $\Gamma$:} integration by parts puts the flux terms at $\tilde x$;
        the physically known quantity (continuous flux) lives at
        $x_\Gamma$. For $f = 0$ the two coincide on each side (the flux of
        the extension is constant in 1D); the extended flux average is
        exactly what cancels the IBP terms. This mirrors the structure of
        the Dirichlet SBM form in the paper, where the first flux term
        also carries the \emph{plain} test trace and only the adjoint
        term carries $\mathcal E$.
  \item \textbf{Adjoint term and penalty on $[\mathcal E u]_\Gamma$:} the
        value-continuity condition holds at $x_\Gamma$; both terms vanish
        on the exact solution, so any consistent-looking location would
        do --- placing them at $\Gamma$ in extended variables enforces
        the correct condition to full order.
  \item \textbf{Shifted source \eqref{eq:sim_rhs}:} with $f \ne 0$ the
        one-sided fluxes at $\tilde x$ and at $x_\Gamma$ differ by the
        strip integral, $\kappa u'(\tilde x) = \kappa u'(x_\Gamma) +
        \int_{\tilde x}^{x_\Gamma} f$; for $f$ continuous this shift is
        the \emph{same on both sides}, hence it multiplies
        $[v]_{\tilde x}$ and moves to the right-hand side. Omitting it
        costs $O(h)$ consistency (observed: convergence stalls). For $f$
        discontinuous across $\Gamma$ the two shifts differ and an
        additional $\{v\}$-weighted data term appears --- same derivation.
\end{itemize}

\subsection*{I.4 Numerical validation (see \texttt{interface\_1d.py})}

$\kappa^- = 1$, $\kappa^+ = 10$, $x_\Gamma = 0.523$ (generic, off-node),
$\mathbb P_p$ elements, $n = 8$ elements unless stated:

\begin{center}
\begin{tabular}{lccc}
\hline
 & $p=1$ & $p=2$ & $p=3$ \\
\hline
broken line ($f=0$), consistent form & $8\cdot10^{-15}$ & $5\cdot10^{-14}$ & $8\cdot10^{-14}$ \\
broken line, naive form              & $3\cdot10^{-2}$ & $8\cdot10^{-2}$ & $2\cdot10^{-1}$ \\
$f=1$ (piecewise quadratic), $p\ge2$ &  --- & $7\cdot10^{-14}$ & $6\cdot10^{-14}$ \\
\hline
\end{tabular}
\end{center}

The consistent form reproduces every exact solution lying in the discrete
space to machine precision (broken line for all $p$; the piecewise
quadratic for $p \ge 2$, which certifies the shifted-source term), while
$p=1$ converges at the expected rate. Robustness in the jump: with
$\kappa^+/\kappa^- = 10^6$ the broken line is reproduced to $10^{-8}$
(conditioning-limited), thanks to the harmonic weights
$\beta^\mp \propto \kappa^\pm$ and the harmonic-mean penalty --- the
standard cure from unfitted Nitsche interface methods, which carries over
unchanged.

\subsection*{I.5 Remarks and outlook to 2D/3D}

\begin{itemize}
  \item \textbf{Relation to the literature.} A shifted interface method
        of this flavor exists (Li \& Scovazzi and follow-ups); these
        notes' added value is the explicit identification of (a) the
        location mismatch $[v]_{\tilde x}$ vs.\ $[\mathcal E v]_\Gamma$ as
        the reason naive derivations fail already in 1D, and (b) the
        shifted-source correction \eqref{eq:sim_rhs}, which is easy to
        miss (it vanishes in every $f=0$ sanity test) yet caps accuracy
        at first order when omitted.
  \item \textbf{No mixed formulation is needed.} Flux continuity is
        enforced through the extended-flux average exactly as value
        continuity is enforced through the extended jump; both are
        Nitsche-type terms on the primal variable.
  \item \textbf{In 2D} the surrogate interface is the staircase of faces
        nearest to $\Gamma$, the strip cells are assigned to one side
        with that side's $\kappa$, and all findings of the definiteness
        notes (theory.tex) apply with two changes in our favor: the
        shift satisfies $|d| \le h/2$ per face by nearest-face selection
        (vs.\ $\sqrt2 h$ for the Dirichlet staircase pockets), and both
        traces are one-sided extrapolations over that reduced distance.
        The same Chebyshev certificates should be run for the coupled
        two-cell block (two fields, duplicated interface DoFs); the
        chamfer and the PDE-substituted extension
        ($\partial_n^2 u = -(f + \nabla\kappa\cdot\nabla u)/\kappa
        - \Delta_{\mathrm{tan}} u$ per side, $\kappa$ constant per side)
        transfer verbatim.
  \item \textbf{Smoother.} The duplicated interface DoFs must belong to a
        common patch (the vertex patch at $\tilde x$ contains both
        one-sided fields by construction), so the full-residual shy patch
        machinery extends without structural change.
  \item \textbf{Open point.} The first flux term in \eqref{eq:sim_form}
        is non-symmetric by construction (plain vs.\ extended test
        traces); the effect on the definiteness certificates at
        $p \ge 3$, and whether the interface analogue prefers the
        PDE-substituted extension from the start, should be checked with
        the same single-block methodology before a 2D implementation.
\end{itemize}
 at the end of the paper to compile.
% ==========================================================================

\section*{Appendix (working notes): a consistent shifted formulation for interface problems}

\subsection*{I.1 Problem and setup}

Model problem on $(0,1)$ (all statements extend verbatim to $d>1$):
\[
  -(\kappa u')' = f, \qquad
  \kappa = \begin{cases} \kappa^- & x < x_\Gamma \\ \kappa^+ & x > x_\Gamma \end{cases},
  \qquad
  [u]_{x_\Gamma} = 0, \quad [\kappa u']_{x_\Gamma} = 0,
\]
with the interface $x_\Gamma$ \emph{not} aligned with the mesh. Shifted
treatment: the surrogate interface is the nearest mesh vertex $\tilde x$,
shift $d = x_\Gamma - \tilde x$ (so $|d| \le h/2$ by construction --- note
this is automatically milder than the Dirichlet SBM staircase, where $s$
reaches $\sqrt 2$). The degree of freedom at $\tilde x$ is duplicated: two
one-sided fields $u_h^\pm$, each assembled with its \emph{own} constant
$\kappa^\pm$ on its own surrogate subdomain.

Interpretation that makes everything work: $u_h^\pm$ do \emph{not}
approximate the physical solution near the interface; they approximate the
smooth \emph{extensions} $\tilde u^\pm$ of the two physical pieces, each
solving its own $\kappa^\pm$-equation past $x_\Gamma$. In the strip between
$\tilde x$ and $x_\Gamma$, the physical solution is recovered in
post-processing from the field of the side that physically owns the strip
(evaluated by extension). The interface conditions are imposed at the
\emph{true} interface via Taylor/polynomial extension
$\mathcal E^\pm w (x_\Gamma) = w(\text{eval.\ of the adjacent element's
polynomial at } x_\Gamma)$, exactly as in the Dirichlet SBM. One side
evaluates by extrapolation (distance $|d|$), the other by interior
evaluation. No mixed formulation, no extra flux variable.

\subsection*{I.2 Why the natural symmetric form fails (diagnosis of the 1D failure)}

Notation: plain jump at the surrogate vertex $[v]_{\tilde x} = v^-(\tilde
x) - v^+(\tilde x)$; extended jump at the true interface $[\mathcal E
v]_\Gamma = \mathcal E^- v^- (x_\Gamma) - \mathcal E^+ v^+(x_\Gamma)$;
weighted average of extended fluxes $\{\kappa \mathcal E u'\}_\Gamma =
\beta^- \kappa^- \mathcal E^- u'\,(x_\Gamma) + \beta^+ \kappa^+ \mathcal
E^+ u'(x_\Gamma)$ with the harmonic weights $\beta^\mp =
\kappa^\pm/(\kappa^-{+}\kappa^+)$.

The trap is a mismatch of \emph{locations}: integration by parts of the
exact solution over the two surrogate subdomains produces flux terms
\emph{at $\tilde x$} acting on the \emph{plain} jump $[v]_{\tilde x}$,
whereas the physical conditions hold \emph{at $x_\Gamma$} in the
\emph{extended} variables. The symmetric-looking Nitsche form that pairs
the flux with $[\mathcal E v]_\Gamma$,
\[
  \cdots - \{\kappa \mathcal E u'\}_\Gamma\, [\mathcal E v]_\Gamma - \cdots
  \qquad \text{(``naive'')},
\]
leaves, for the exact broken line with flux $F$ (take $f=0$, $p=1$), the
residual
\[
  F\big( [v]_{\tilde x} - [\mathcal E v]_\Gamma \big)
  = -F\, d\, [v']_{\tilde x} \;\neq\; 0 ,
\]
an $O(d) = O(h)$ \emph{consistency} error concentrated at the interface:
the method converges to the wrong solution, and already the piecewise
linear exact solution is not reproduced. This is precisely the observed 1D
failure. (Verified numerically: the naive form gives errors
$3\cdot10^{-2}$--$2\cdot10^{-1}$ on the broken line, non-decreasing in
$p$.)

\subsection*{I.3 The consistent form}

The two interface conditions play different roles and must live at
different locations:

\begin{equation}
  a(u,v) \;=\; \sum_\pm \int_{\tilde\Omega^\pm} \kappa\, u' v'
  \;-\; \{\kappa \mathcal E u'\}_\Gamma\, \underbrace{[v]_{\tilde x}}_{\text{plain, at }\tilde x}
  \;-\; \{\kappa \mathcal E v'\}_\Gamma\, \underbrace{[\mathcal E u]_\Gamma}_{\text{extended, at }\Gamma}
  \;+\; \sigma\, [\mathcal E u]_\Gamma [\mathcal E v]_\Gamma ,
  \label{eq:sim_form}
\end{equation}
with $\sigma = \gamma\, \kappa_{\mathrm{hm}}\, p^2/h$,
$\kappa_{\mathrm{hm}} = 2\kappa^-\kappa^+/(\kappa^-{+}\kappa^+)$, and the
right-hand side carrying a \emph{shifted-source correction}:
\begin{equation}
  L(v) \;=\; \sum_\pm \int_{\tilde\Omega^\pm} f v
  \;+\; \Big( \int_{\tilde x}^{x_\Gamma} f \, dx \Big)\, [v]_{\tilde x} .
  \label{eq:sim_rhs}
\end{equation}

Rationale, term by term:
\begin{itemize}
  \item \textbf{First flux term at $\tilde x$, extended flux from
        $\Gamma$:} integration by parts puts the flux terms at $\tilde x$;
        the physically known quantity (continuous flux) lives at
        $x_\Gamma$. For $f = 0$ the two coincide on each side (the flux of
        the extension is constant in 1D); the extended flux average is
        exactly what cancels the IBP terms. This mirrors the structure of
        the Dirichlet SBM form in the paper, where the first flux term
        also carries the \emph{plain} test trace and only the adjoint
        term carries $\mathcal E$.
  \item \textbf{Adjoint term and penalty on $[\mathcal E u]_\Gamma$:} the
        value-continuity condition holds at $x_\Gamma$; both terms vanish
        on the exact solution, so any consistent-looking location would
        do --- placing them at $\Gamma$ in extended variables enforces
        the correct condition to full order.
  \item \textbf{Shifted source \eqref{eq:sim_rhs}:} with $f \ne 0$ the
        one-sided fluxes at $\tilde x$ and at $x_\Gamma$ differ by the
        strip integral, $\kappa u'(\tilde x) = \kappa u'(x_\Gamma) +
        \int_{\tilde x}^{x_\Gamma} f$; for $f$ continuous this shift is
        the \emph{same on both sides}, hence it multiplies
        $[v]_{\tilde x}$ and moves to the right-hand side. Omitting it
        costs $O(h)$ consistency (observed: convergence stalls). For $f$
        discontinuous across $\Gamma$ the two shifts differ and an
        additional $\{v\}$-weighted data term appears --- same derivation.
\end{itemize}

\subsection*{I.4 Numerical validation (see \texttt{interface\_1d.py})}

$\kappa^- = 1$, $\kappa^+ = 10$, $x_\Gamma = 0.523$ (generic, off-node),
$\mathbb P_p$ elements, $n = 8$ elements unless stated:

\begin{center}
\begin{tabular}{lccc}
\hline
 & $p=1$ & $p=2$ & $p=3$ \\
\hline
broken line ($f=0$), consistent form & $8\cdot10^{-15}$ & $5\cdot10^{-14}$ & $8\cdot10^{-14}$ \\
broken line, naive form              & $3\cdot10^{-2}$ & $8\cdot10^{-2}$ & $2\cdot10^{-1}$ \\
$f=1$ (piecewise quadratic), $p\ge2$ &  --- & $7\cdot10^{-14}$ & $6\cdot10^{-14}$ \\
\hline
\end{tabular}
\end{center}

The consistent form reproduces every exact solution lying in the discrete
space to machine precision (broken line for all $p$; the piecewise
quadratic for $p \ge 2$, which certifies the shifted-source term), while
$p=1$ converges at the expected rate. Robustness in the jump: with
$\kappa^+/\kappa^- = 10^6$ the broken line is reproduced to $10^{-8}$
(conditioning-limited), thanks to the harmonic weights
$\beta^\mp \propto \kappa^\pm$ and the harmonic-mean penalty --- the
standard cure from unfitted Nitsche interface methods, which carries over
unchanged.

\subsection*{I.5 Remarks and outlook to 2D/3D}

\begin{itemize}
  \item \textbf{Relation to the literature.} A shifted interface method
        of this flavor exists (Li \& Scovazzi and follow-ups); these
        notes' added value is the explicit identification of (a) the
        location mismatch $[v]_{\tilde x}$ vs.\ $[\mathcal E v]_\Gamma$ as
        the reason naive derivations fail already in 1D, and (b) the
        shifted-source correction \eqref{eq:sim_rhs}, which is easy to
        miss (it vanishes in every $f=0$ sanity test) yet caps accuracy
        at first order when omitted.
  \item \textbf{No mixed formulation is needed.} Flux continuity is
        enforced through the extended-flux average exactly as value
        continuity is enforced through the extended jump; both are
        Nitsche-type terms on the primal variable.
  \item \textbf{In 2D} the surrogate interface is the staircase of faces
        nearest to $\Gamma$, the strip cells are assigned to one side
        with that side's $\kappa$, and all findings of the definiteness
        notes (theory.tex) apply with two changes in our favor: the
        shift satisfies $|d| \le h/2$ per face by nearest-face selection
        (vs.\ $\sqrt2 h$ for the Dirichlet staircase pockets), and both
        traces are one-sided extrapolations over that reduced distance.
        The same Chebyshev certificates should be run for the coupled
        two-cell block (two fields, duplicated interface DoFs); the
        chamfer and the PDE-substituted extension
        ($\partial_n^2 u = -(f + \nabla\kappa\cdot\nabla u)/\kappa
        - \Delta_{\mathrm{tan}} u$ per side, $\kappa$ constant per side)
        transfer verbatim.
  \item \textbf{Smoother.} The duplicated interface DoFs must belong to a
        common patch (the vertex patch at $\tilde x$ contains both
        one-sided fields by construction), so the full-residual shy patch
        machinery extends without structural change.
  \item \textbf{Open point.} The first flux term in \eqref{eq:sim_form}
        is non-symmetric by construction (plain vs.\ extended test
        traces); the effect on the definiteness certificates at
        $p \ge 3$, and whether the interface analogue prefers the
        PDE-substituted extension from the start, should be checked with
        the same single-block methodology before a 2D implementation.
\end{itemize}
